\section{Open Problems and Research Agenda}
\label{sec:openproblems}

This section consolidates the open problems flagged
throughout the survey into a single agenda for the next
standardisation cycle. Closing on an agenda of open problems
is an established form for surveys~\cite{parkerholder2022autorl}. The six security-specific items
identified in \S\ref{sec:security:open} (O1 to O6) carry
forward into this section as the substrate for the
broader research agenda; we re-label them here as items
P1 to P6 to mark that the open-problems agenda is broader
than security alone. P7 to P12 are problems flagged
elsewhere in the survey: the Axis 9 schema gap, the
cross-network mandate harmonisation, the trust-bootstrap
problem across heterogeneous registries, the
benchmarking gap, the economic-incentive question, and
the regulation-readiness gap. The dozen items together
constitute what the survey treats as the agenda the
field should address by the end of the 2026 to 2027
revision cycle.

Section~\ref{sec:openproblems:convergence} treats the
convergence-versus-fragmentation question that frames the
rest of the agenda. Sections~\ref{sec:openproblems:trust}
through~\ref{sec:openproblems:regulation} unpack five
problem clusters, each grouping a small number of P-items
that share a common substrate. The survey takes positions
where the literature is clear and refrains from positions
where the literature is open; the closing
section~\ref{sec:openproblems:priorities} ranks the
P-items by tractability and impact for the benefit of
researchers selecting where to work.

\subsection{Convergence versus fragmentation}
\label{sec:openproblems:convergence}

The framing question for the entire open-problems agenda
is whether the contemporary protocols converge to a
single dominant pattern or remain a co-stack of
specialised protocols indefinitely. The survey's
position, after weighing the evidence from
Sections~\ref{sec:industry-protocols} to \ref{sec:governance},
is that \emph{partial convergence at the primitive level
is likely; full convergence at the protocol level is not}.

\paragraph{The convergence argument.}
Three of the four large peer-messaging protocols (A2A,
ANP, AGNTCY) are converging on shared primitives at the
identity, discovery, and task-lifecycle layers under
the AAIF and W3C umbrellas. The agentic commerce stack
is converging on the intent-mandate-receipt envelope
across AP2, Visa TAP, Mastercard Agent Pay, and
ERC-8004. The tool-binding layer has effectively
converged on MCP at the production cutoff. The
direction is unmistakable: the field is moving toward
a smaller set of shared primitives.

\paragraph{The fragmentation counter-argument.}
At the same time, the three architectural layers of
\S\ref{sec:taxonomy:layers} are not collapsing into a
single layer. A2A and MCP remain structurally distinct
because the questions they answer are structurally
distinct; the commerce-settlement layer remains
structurally distinct from peer-messaging because the
intent-mandate-receipt envelope has no analogue in
inter-agent task semantics. The pattern is convergence
\emph{within} layers and persistent diversity
\emph{across} layers. This is consistent with how
HTTP, TLS, and DNS converged within their respective
layers while remaining different protocols.

\paragraph{Implication for the research agenda.}
The convergence-versus-fragmentation framing affects
what research direction the field should reward.
Within-layer convergence rewards primitive-level
research (better identity primitives, better
delegation primitives, better content models)
because the primitives will be adopted across
multiple protocols. Cross-layer divergence rewards
composition-pattern research (how layers stack
cleanly, what bridges between them look like, how
end-to-end traces flow across layers) because that
work is not subsumed by any single protocol's
specification.

\subsection{Trust bootstrapping and identity portability}
\label{sec:openproblems:trust}

The first cluster of P-items concerns identity and the
trust bootstrap that supports it.

\paragraph{P1: Cryptographic delegation chains.}
The single most-cited gap (O1 of
\S\ref{sec:security:open}). No surveyed protocol
standardises a delegation credential, and token forwarding
remains the de-facto pattern deployments adopt in its
absence. The candidate successor patterns are the Agentic
JWT~\cite{agentic_jwt} and the
authorisation-propagation chain~\cite{authz_propagation}.
ACNBP's capability-binding
ceremony~\cite{spec_acnbp} is a third ingredient rather
than a third candidate: it makes one hop verifiable
without an authorisation server, and it says nothing about
composing hops. The research question is which
of the candidates the field standardises and how the
chosen pattern bridges the OAuth-bearer-token
deployments to a cryptographic-chain world without
breaking deployed implementations. The AGRP envelope
(Contribution C6, \S\ref{sec:proposals:agrp}) provides a
container in which any of these patterns can be carried
forward as the \texttt{filter\_signature} chain; AGRP does
not specify the chain mechanism, only the envelope that
transports it.

\paragraph{P2: Identity rotation.}
O2 of \S\ref{sec:security:open}. No contemporary
protocol specifies how an agent's identifier rotates
without breaking in-flight tasks. The research question
is whether the rotation primitive belongs at the
identity layer (DID document updates with grace
periods) or at the task layer (Task lifecycle
extensions for identifier migration); the survey
takes no position.

\paragraph{P7: Trust bootstrapping across heterogeneous
registries.}
A new problem not previously enumerated. Agents
registered under different discovery infrastructures
(A2A AgentCards under one well-known endpoint, ANP
DIDs under a DID resolver, Coral agents under a
Coral Server, AGNTCY agents in an OASF directory) must
trust each other across the boundary. No protocol
specifies how this cross-registry trust bootstraps:
how does an A2A-registered agent verify the
authenticity of an ANP-registered counterpart?
The naive answer (each registry must be trusted by
each consumer) does not scale; the candidate
mechanisms (cross-registry attestation, federated
trust anchors, blockchain-anchored cross-resolution)
are at the pre-specification stage. The OWA scoping
scheme (Contribution C5, \S\ref{sec:proposals:owa})
partially addresses this by anchoring organizational
trust in DNSSEC and providing a normative workspace
tier across which trust can be expressed; the
cross-registry attestation question remains open.

\subsection{Cross-protocol bridging and gateways}
\label{sec:openproblems:bridging}

The second cluster concerns the composition and bridging
mechanisms between layers and between protocols within a
layer.

\paragraph{P8: ANP-to-A2A identity bridges.}
The composition that the rubric matrix most clearly
recommends (\S\ref{sec:comparison:stack}: ANP's
identity-layer strengths + A2A's task-layer strengths)
is not yet a deployed stack. The bridge specification
would say: an ANP-resolved DID presents itself across
an A2A AgentCard with the AgentCard fields populated
from the ANP AgentDescription. The AAIF working group
has begun the bridge but has not published a
specification at the cutoff. This is the single
composition that the survey's rubric most clearly
indicates would deliver more capability per spec-word
than any other.

The scope of the bridge is narrower than the protocol pair,
and stating it narrowly is what makes the problem
tractable. ANP is a suite of sub-protocols, and only the
identity mechanism composes cleanly: a DID carries no
dependency on the rest of ANP. Agent Description and Agent
Discovery overlap what the A2A AgentCard already provides,
so the bridge must decide which descriptor is
authoritative rather than populate one from the other in
both directions (\S\ref{sec:comparison:stack}). The
tractable version of P8 is a one-way binding: the DID
authenticates the agent, the AgentCard describes it, and
the bridge specifies only how the card commits to the DID.
Attempting to reconcile two descriptor formats is a
different and larger problem, and conflating the two is
one reason the bridge has not shipped.

\paragraph{P9: Cross-network mandate harmonisation in
the commerce stack.}
The commerce-settlement layer leaves a cross-network
interoperability gap (\S\ref{sec:emerging:commerce},
\S\ref{sec:adoption:commerce}): a single agent
mandate should be portable across PayPal AP2, Visa
TAP, Mastercard Agent Pay, and ERC-8004 settlement
substrates without re-signing per network. The AAIF
Agentic Commerce Working Group~\cite{lf_aaif} has
begun the harmonisation; the open question is whether
the harmonised envelope is a wire-format unification
(all four protocols carry the same envelope schema)
or a semantic unification (different schemas with a
specified translation) or a higher-level meta-mandate
that wraps the four substrates.

\paragraph{P10: The gateway-versus-translator
question.}
A meta-level open question raised by P8 and P9: should
the bridges between protocols be \emph{gateways} (a
deployed intermediary that translates between
protocols at the message boundary) or \emph{translators}
(a specification primitive that each protocol's
implementations consume directly)? The two answers
have different governance, deployment, and security
implications. The SAGA architecture~\cite{saga2026}
implicitly chooses the gateway pattern; the OASF
descriptor~\cite{spec_agntcy} implicitly chooses the
translator pattern. The field has not yet decided.

\subsection{Benchmarking gaps}
\label{sec:openproblems:benchmarking}

The third cluster concerns the evaluation
infrastructure that would let the field measure
progress on the other open problems.

\paragraph{P11: A rubric-stable benchmark for
inter-agent protocols.}
The Axis 8 security benchmark
(A2ASecBench~\cite{a2asecbench}) is the most mature
empirical benchmark at the cutoff, but it covers
only one rubric axis and only one protocol family
(A2A). MultiAgentBench~\cite{multiagentbench} and
COMMA~\cite{comma_tmlr25} cover the collaboration
and multimodal-coordination dimensions but do not
align with the rubric axes.
AgentsNet~\cite{agentsnet} measures coordination and
collaborative reasoning over agent graphs, and
H\"arer~\cite{spec_eval_cyber} evaluates specified
multi-agent systems on cybersecurity tasks. Both
measure agent behaviour rather than what a
specification requires, which is the distinction the
rubric turns on. A rubric-stable
benchmark would cover all ten axes and would let
the field measure protocol revisions against a
fixed scoring instrument. The ``Which LLM
Multi-Agent Protocol to Choose''
benchmark~\cite{protocolbench} is the closest
candidate to this aim at the cutoff but is
implementation-leaning rather than spec-leaning.
The research question is whether a spec-leaning
benchmark of the rubric's form is practical and
who maintains it.

A complementary implementation track should preserve the
rubric score while measuring behavior. For privacy, this
requires capture-completeness indicators and results by
actual disclosure channel, including at minimum inter-agent
messages, shared memory, tool-call arguments, and final
output, rather than a single output-only pass rate.
Source channels (user input and tool response) and
lifecycle stages (pool construction, selection, execution,
and output) must be represented separately. A benchmark
that evaluates context selection should also join each
content-free selection decision to later disclosures by a
stable candidate reference and retain logging or fallback
failures as explicit observations. Risk-weighted summaries
may supplement these primary measurements, but an
unpublished scoring method cannot serve as a citable
primary outcome.

Two prerequisites of that benchmark are missing and are
worth naming as work items rather than assumptions. The
first is a two-rater pass over the rubric. This survey
reports calibration rules, worked examples, and an auditable
change history, and no independent second scoring
(\S\ref{sec:rubric:calibration}). A benchmark built on a
scoring instrument whose reproducibility has never been
measured inherits that gap. The axes to measure first
are the two most judgement-heavy ones: semantic
interoperability (Axis 9) and the observability component of
Axis 10. The
second is a labelled corpus of inter-agent messages with
ground-truth sensitive spans. Without it, no egress filter's
false-negative and false-positive rates can be reported, so
no privacy-preserving envelope can be evaluated on the
property that actually matters
(\S\ref{sec:proposals:agrp:open}). Both are cheaper than the
benchmark they enable, and neither exists at the cutoff.

\paragraph{P12: Multi-tenant evaluation harness.}
Existing evaluation work assumes single-tenant
deployments. The cross-tenant policy enforcement
question (P4) requires an evaluation harness that
simulates multiple tenants with conflicting
policies and measures whether the protocol's
implementations enforce isolation correctly. No
such harness exists at the cutoff. The research
question is what the harness should require of
the protocols and whether the harness lives at the
protocol-conformance layer or the deployment
layer.

\paragraph{Benchmarking discipline beyond AI evals.}
A meta-observation that the survey should record:
the inter-agent benchmark space is structurally
different from the LLM-evaluation space because
the unit of comparison is a specification (or an
implementation conformance) rather than a model. The
existing benchmark discipline does not import cleanly,
and the open-problems agenda for this area is more
methodological than the others: the field needs to
develop benchmarking conventions appropriate to
protocol specifications.

\subsection{Economic and incentive design}
\label{sec:openproblems:economic}

The fourth cluster concerns the economic incentives
that will shape adoption.

\paragraph{Pricing and metering.}
A protocol-level question that the survey has
not yet addressed: how are agent-invoked services
priced and metered? The dominant pattern at the
cutoff is per-invocation pricing
(every MCP tool call is a billable event), but
the commerce-settlement layer enables mandate-bounded
pricing (the principal authorises spending up to
budget X for outcome Y, regardless of how many
invocations it takes). The two pricing models
have different governance and audit implications;
no contemporary protocol specifies which model the
protocol assumes, leaving the question to
deployment.

\paragraph{Reputation and ranking.}
A second open question: how does the discovery
layer rank competing agents? The naive answer is
that the registry returns all matches and the
caller chooses; the practical answer is that
discovery layers will increasingly include
reputation signals (success rate, response time,
audit cleanliness) that influence ranking. The
research question is whether the protocol
specifies the reputation primitive or leaves it
to the discovery infrastructure, and how the
reputation signal resists manipulation.

\paragraph{Network-effect protection.}
A third open question, motivated by
\S\ref{sec:governance}: how does the
inter-agent protocol ecosystem avoid the
winner-take-all dynamic that the network-effect
literature predicts for similar-shaped
infrastructures? The
license-permissiveness argument
(\S\ref{sec:governance:foundations}) is one
mechanism; the AAIF founding-member breadth is
another; the open question is whether these
mechanisms are sufficient or whether
additional structural protection is needed.

\subsection{Regulation-ready protocol design}
\label{sec:openproblems:regulation}

The fifth cluster concerns the protocols' alignment
with the regulatory environment they will operate in.

\paragraph{P3: Prompt-injection mitigation as a
spec-level requirement.}
O3 of \S\ref{sec:security:open}. Regulators in
multiple jurisdictions are converging on the
position that agent operators are responsible
for protecting principals against adversarial
content; the protocol-level question is
whether the protocols themselves carry the
mitigation requirement or leave it to
operator-level policy. ACNBP and LOKA are the
only surveyed protocols that include
prompt-injection mitigation as part of the
specification; the research question is whether
the AAIF working groups will follow. The AGRP
envelope (Contribution C6,
\S\ref{sec:proposals:agrp}) makes
prompt-injection mitigation a spec-level requirement
by mandating that the egress filter run before the
message is signed; AGRP partially addresses P3 by
making the requirement normative and verifiable,
though the filter rules themselves are
organization-controlled rather than
protocol-mandated.

\paragraph{P4: Cross-tenant policy enforcement.}
O4 of \S\ref{sec:security:open}. Multi-tenant
agentic deployments need per-tenant policy
isolation that the contemporary protocols do
not provide. The SAGA architecture is the
leading proposal; the research question
includes whether SAGA's TEE-anchored gateway
becomes a normative requirement, an
implementation pattern, or a regulatory
recommendation.

\paragraph{P5: Contextual integrity for
forwarding.}
O5 of \S\ref{sec:security:open}. The
conversation-history forwarding pattern leaks
content that the principal did not authorise
the remote agent to receive. Regulatory
attention on this dimension is increasing
(particularly in healthcare and finance
verticals where forwarding crosses regulated
boundaries). The protocol-level question is
whether the forwarding decision is made by the
protocol (with a specified contextual-integrity
primitive) or by the orchestrator (with the
protocol providing only the metadata to enable
the decision). The
\emph{Beyond Message Passing} position
\cite{beyond_message_passing} argues for the
protocol-layer answer; the
\emph{Semantic-Driven} position
\cite{semantic_driven} argues for the
orchestrator-layer answer. The survey proposes
ACSP (Contribution C7,
\S\ref{sec:proposals:acsp}) as the protocol-layer
candidate: a wire-level selection envelope that
records which items the sender forwarded, which
it dropped, and an attestation chain the receiver
can verify. AGRP's classification labels
(Contribution C6,
\S\ref{sec:proposals:agrp:schema}) provide the
metadata an orchestrator-layer decision would
also need, and ACSP composes with AGRP rather
than replacing it
(\S\ref{sec:proposals:acsp:composition}). The
orchestrator-versus-protocol dispatch remains
open: whether ACSP displaces the orchestrator-layer
alternatives or composes with them depends on
empirical evidence the field does not yet have.

The missing empirical design has two parts. First, the
same candidate must be evaluated under declared
same-workspace, cross-workspace, and cross-organisation
boundaries so that authorisation is not inferred from
semantic relevance. Second, a content-free selection
record must be joined to later observations on the actual
inter-agent, shared-memory, tool-call, and final-output
channels. Hard policy violations, graded risk, necessary
context recall, and end-task success should be reported as
separate outcomes; an unavailable channel must not be
counted as a non-disclosure.

\paragraph{P6: Side-channel metadata privacy.}
O6 of \S\ref{sec:security:open}. The metadata
associated with inter-agent messages leaks
information about the principal's actions even
when the content is encrypted. No contemporary
protocol addresses action privacy in the
metadata layer. The research question is
whether action-privacy primitives belong in the
protocol (ACNBP's ring signatures suggest one
direction) or in a privacy-preserving deployment
layer. AGRP (Contribution C6) provides
metadata-field encryption hooks within the
envelope but does not specify the
metadata-privacy primitive normatively; this
remains an open question for a future revision
of AGRP.

\subsection{The schema fragmentation problem}
\label{sec:openproblems:schema}

A standalone problem that the survey treats as the
single largest open question.

\paragraph{P-axis-9: Escaping the level-1
semantic-interoperability equilibrium.}
The Axis 9 gap of \S\ref{sec:semantic}. No
contemporary protocol ships a level-3 mechanism
for capability semantics; the field is unsure
that closing the gap is the right objective. The
research question is the one
\S\ref{sec:semantic:fragmentation} closes with:
ontology-base or semantic-matchmaker, and at
what layer (protocol or orchestrator) does the
mechanism live. The
\emph{Beyond Message Passing} and
\emph{Semantic-Driven} lines disagree; the
consequences (the substantially lower
inter-organisational call success rate
documented by the AAIF interoperability working
group~\cite{lf_aaif}, the retrofit cost when
adding schemas after deployment, and the
matchmaker drift across foundation-model
upgrades) are visible enough that the gap will
be addressed within the review cycle. The form
the address takes is not yet legible, though one piece of
post-cutoff evidence points away from the ontology-base
answer: ANP, the only protocol shipping a level-2
mechanism, narrowed its JSON-LD commitment at the
description layer in its 1.1 release and cited
readability~\cite{spec_anp_v11}. A researcher choosing
between the two families should treat ontology adoption as
reversible rather than monotonic
(\S\ref{sec:semantic:ontology}). The OWA workspace tier
(Contribution C5, \S\ref{sec:proposals:owa})
narrows the cross-organizational schema gap by
making per-workspace policy enforcement
explicit, but does not specify the capability
ontology or matchmaker; the Axis 9 gap remains
the largest unresolved research-and-standardization
question even after the C5/C6/C7 triplet.

\subsection{Priorities for the next revision cycle}
\label{sec:openproblems:priorities}

The agenda is large. The survey closes with a
priority ranking organised by tractability and
impact, for the benefit of researchers selecting
where to work.

\paragraph{Tier 1: High impact, near-term
tractable.}
The cryptographic delegation chain (P1) is the
single highest-priority item, and it is the one item on
this agenda that no specification currently addresses at
all. The ingredients are well-developed (the Agentic JWT
derivation pattern, ACNBP's single-hop capability
binding), the deployment urgency is high (token forwarding
is in production today), and the AAIF working group
is the natural institutional home for the
specification. A2A's extension mechanism is the
plausible route: it lets a candidate chain be
developed and evaluated in the field before any
proposal to move it into the core specification. The route
is Microsoft's suggestion, made when we contacted it about
this survey, rather than ours. The ANP-to-A2A identity
bridge (P8) is the second-highest priority for
similar reasons.

\paragraph{Tier 2: High impact, multi-cycle.}
The schema-fragmentation problem
(P-axis-9), cross-network mandate harmonisation
(P9), and cross-tenant policy enforcement (P4)
are high-impact but require multi-cycle work.
The schema problem requires either a new
vocabulary or a matcher architecture, both of
which need substantial development before
specification; the mandate harmonisation
requires multi-network coordination at the
governance level; the cross-tenant problem
requires both protocol specification and
deployment-infrastructure work.

\paragraph{Tier 3: Important, longer-term.}
Prompt-injection mitigation as a spec
requirement (P3), contextual integrity for
forwarding (P5), and identity rotation (P2) are
important problems that the field is not yet
ready to address normatively. The candidate
mechanisms are not mature enough; the research
work to mature them is the prerequisite for the
specification work.

\paragraph{Tier 4: Methodological.}
The benchmarking gap (P11, P12) is at a
different level than the other items: it is
the instrument the field needs to measure
progress on the other items. Tier-4 in
priority is not Tier-4 in importance; the
benchmark work is necessary for the rest of
the agenda to be tractable. The cheapest item in
this tier applies to this survey's own instrument: a
second rater scoring Axis 9 and the observability
component of Axis 10, the two most judgement-heavy
axes. This survey reports no inter-rater agreement
statistic (\S\ref{sec:rubric:calibration}), so the
reproducibility of the rubric is at present an
argued property rather than a measured one.

\paragraph{What this section does not predict.}
The survey closes the open-problems agenda
with the caveat that the empirical
deployment evidence (\S\ref{sec:adoption})
shows adoption density uncorrelated with
technical rigour. The agenda above lists
what the survey considers the technically
important problems; the field will work on
what the institutional incentives reward. The
two sets overlap but do not coincide, and the
survey's role is to make the technical agenda
legible enough that researchers and
practitioners can advocate for the
high-impact items even when the institutional
incentives point elsewhere.
